\section{Tensor class}
\label{sec:autodiff}

\subsection{Definition}
A \code{Tensor} holds an array \code{data} (always \code{float64}), a flag \code{requires\_grad},
and, when the flag is set, three things about how it was made: the set of tensors it was computed
from (\code{\_prev}), a closure \code{\_backward}, and the name of the operation. Every operation
on tensors returns a new tensor, and if any input requires a gradient it records those three
fields. A forward pass therefore builds a directed acyclic graph whose leaves are the inputs and
the parameters, and whose root is, in training, the scalar loss $L$.

\subsection{The backward pass}
For a tensor $u$ in the graph, $\dL{u}$ is the gradient of the loss with respect to it: an array
with the same shape as $u$, stored in \code{u.grad}. If $v$ was computed from $u$ and $\dL{v}$ is
already known, the chain rule gives
\begin{equation}
  \frac{\partial L}{\partial u} = J\T\, \frac{\partial L}{\partial v},
  \qquad J = \frac{\partial v}{\partial u},
  \label{eq:chain}
\end{equation}
with $u$ and $v$ read as flat vectors and $J$ the Jacobian of $v$ with respect to $u$. This product
of a vector with the Jacobian is what each operation's \code{\_backward} computes for each of its
inputs, without ever building $J$: for an elementwise operation it is diagonal, and the product is
just another elementwise operation on arrays.

A tensor used by more than one operation receives a gradient from each, and by the chain rule they
add up. That is why every \code{\_backward} does \code{u.grad += \dots} and never \code{u.grad = \dots}.

\code{Tensor.backward()} first orders the graph topologically with an iterative depth-first
search, then sets $\partial L/\partial L = 1$ and calls \code{\_backward} on every node in reverse
topological order, so that each node's gradient is complete before it is passed on to its inputs.

\begin{algorithm}[H]
\caption{\code{Tensor.backward} (reverse mode)}
\begin{algorithmic}[1]
\State $\textit{order} \gets$ topological order of the graph reachable from $L$
\State \code{L.grad} $\gets 1$
\For{$v$ in reverse(\textit{order})}
  \State $v.\code{\_backward}()$ \Comment{adds its part of the gradient to the \code{.grad} of each input}
\EndFor
\end{algorithmic}
\end{algorithm}

Gradients are allocated lazily (as zeros, the first time they are read) and must be reset between
steps with \code{zero\_grad()}, which every optimizer does before evaluating the loss.

\subsection{Broadcasting}
NumPy broadcasting lets $a \in \R^{n\times m}$ be added to $b \in \R^{m}$. The forward is free; the
backward is not, because \code{b.grad} must have the shape of $b$. Since $b$ was copied along the
broadcast axes, its gradient is the sum of the incoming gradient over them. \code{reduce\_grad}
implements exactly this: it sums over the leading axes that broadcasting added, and over every axis
where $b$ had size $1$ but the result did not.

\subsection{Freezing a model}
A trained model used only for predictions does not need a graph: building one keeps every
intermediate array alive for a backward pass that never comes. \code{Module.freeze()} sets
\code{requires\_grad = False} on every parameter, after which a forward pass on a plain input
records nothing. A gradient with respect to an input that does require one is still computed
exactly, which is all that a saliency map needs.
